\section{Objetivos}
El objetivo principal de esta práctica de laboratorio es relevar y determinar las curvas características intensidad-tensión ($I$ vs $V$) para distintos componentes eléctricos de dos terminales, con el fin de analizar su linealidad.

Los elementos específicos a evaluar son:
\begin{itemize}
    \item Resistencia de carbón depositado.
    \item Lámpara de filamento de tungsteno.
    \item Diodo semiconductor de silicio.
    \item Diodo emisor de luz (LED).
\end{itemize}

\vspace{0.5cm}

\section{Resumen}
Para comprender la naturaleza de la conducción eléctrica en distintos dispositivos, resulta fundamental caracterizar empíricamente la relación entre la diferencia de potencial aplicada a un componente y la corriente que lo atraviesa. Esta experiencia tuvo como propósito contrastar el comportamiento de elementos óhmicos (lineales) frente a elementos no óhmicos (no lineales), identificando empíricamente las desviaciones de la Ley de Ohm provocadas por efectos térmicos o por la naturaleza de los materiales semiconductores.

Para llevar a cabo este análisis, se implementó un circuito de medición compuesto por una fuente de tensión continua variable, un amperímetro dispuesto en serie, un voltímetro en paralelo y el elemento bajo ensayo. Aumentando progresivamente la tensión suministrada — y asegurando no exceder los valores máximos de tensión y corriente admitidos por cada componente — se relevaron los distintos pares de valores de corriente y tensión. Finalmente, con los datos obtenidos, se construyeron las gráficas correspondientes a las curvas características $I=f(V)$ de cada elemento, permitiendo así evaluar visual y analíticamente su comportamiento y linealidad.

\vspace{0.8cm}
